\paragraph{AQ-136-MULTIWEIGHT: independent coefficient disks in every large weight are still not admissible.}
Fix $N_0\geq1$.  One might try to remove the coefficient-tail
obstruction of AQ-136-INTEGRAL-TAIL by imposing all relations
$\sum_jx_{N,j}=T^{[N]}$ for every $N\geq N_0$, using a
separate divided-power simplex disk for each relation.
This construction has the required vertical cohomology,
but its zero-form row has nonzero horizontal $H_1$ in
arbitrarily large weights, already after extension to
$\mathbb Q$.  It therefore fails condition (136.4), both
over $\mathbb Q$ and over $\mathbb Z$.

\begin{proof}
First define the integral construction precisely.  Put
\[
 R=\Gamma_{\mathbb Z}(U_N:N\geq N_0),\qquad
 S=\Gamma_{\mathbb Z}(T),\qquad
 U_N\longmapsto T^{[N]},\quad \operatorname{wt}(U_N)=N.
                                                               \tag{136.32}
\]
In simplicial degree $n$, let
\[
 A_n=\Gamma_{\mathbb Z}(x_{N,j}:N\geq N_0,\ 0\leq j\leq n)
       \otimes\Lambda(e_{N,j}:N\geq N_0,\ 0\leq j\leq n)
       /(\textstyle\sum_j e_{N,j}:N\geq N_0),                \tag{136.33}
\]
where $d x_{N,j}=e_{N,j}$, $d e_{N,j}=0$, both generators
have weight $N$, and their cochain degrees are zero and
one, respectively.  The coefficient map sends
$U_N$ to $\sum_jx_{N,j}$.  Simplicial operators act on
each $N$-family by the sum-over-fibres rule of (136.5).
All divided powers and differential identities are those
of the integral model AQ-136-DP-MODEL; tensoring the
models and taking the filtered union over finite sets
of indices gives the displayed strict PD algebra.
Define
\[
                         B^{\rm multi}=S\otimes_R A.        \tag{136.34}
\]
Equivalently this is the strict PD quotient with
$\sum_jx_{N,j}=T^{[N]}$ and $\sum_j e_{N,j}=0$ for every
$N\geq N_0$.  Eliminating $x_{N,0}$ and $e_{N,0}$
identifies each simplex with the free divided-power
polynomial de Rham algebra over $S$ on the remaining
variables.  In a fixed weight only finitely many $N$
can occur, and in a fixed simplex there are finitely
many variables for each such $N$.  The weight-preserving
integral disk contractions from (136.6) therefore apply
in each weight, proving vertical cohomology $S$ in
degree zero and zero in every positive degree.

We compute the zero-form horizontal homology.  An
invertible linear change of divided-power variables
replaces $x_{N,0}$ by $U_N=\sum_jx_{N,j}$, independently
for every $N$.  Hence $A_n^0$ is a free $R$-module,
with divided-power monomials in $x_{N,j}$, $j\geq1$,
as basis.  Its augmentation to $\mathbb Z$ sends every
variable to zero.  As a simplicial abelian group, its
weight-$w$ part is the direct sum over the finitely
many tuples $(a_N)$ satisfying $\sum_N N a_N=w$ of
\[
                         \bigotimes_N\Gamma^{a_N}(L).
                                                               \tag{136.35}
\]
For $w>0$, at least one $a_N$ is positive.  That factor
is linearly simplicially contractible by the homotopy
calculation preceding (136.5), and tensoring its
homotopy with the identities contracts the displayed
tensor product.  Weight zero is the constant module
$\mathbb Z$.  Thus the augmented normalized complex
$N(A_*^0)\to\mathbb Z$ is exact.  Each normalized term
is an $R$-linear direct summand of a free term $A_n^0$
by the standard degeneracy decomposition.  It is
therefore a projective $R$-module, so this is a
projective resolution of the augmentation module
$\mathbb Z$ over $R$.

Normalization commutes with arbitrary extension of
scalars here: its degeneracy decomposition is split
and functorial, so tensoring carries the projector
onto the normalized summand to the corresponding
projector after tensoring.  Consequently
\[
 H_i\bigl(N((B^{\rm multi})_*^0)\bigr)
                  =\operatorname{Tor}_i^R(S,\mathbb Z).
                                                               \tag{136.36}
\]
Rationalization is flat.  In characteristic zero the
divided-power polynomial algebra is the polynomial
algebra, with $T^{[N]}=t^N/N!$.  Thus (136.36) after
tensoring with $\mathbb Q$ is computed over
\[
 R_{\mathbb Q}=\mathbb Q[u_N:N\geq N_0],\qquad
 S_{\mathbb Q}=\mathbb Q[t],\qquad
                     u_N\longmapsto a_N=t^N/N!.
                                                               \tag{136.37}
\]
The Koszul complex on the countable polynomial
variables $u_N$ is a free resolution of the
augmentation module $\mathbb Q$.  For completeness,
the Koszul complex on a finite initial set is exact
in positive degrees because adjoining a new
variable tensors the previous resolution with the
two-term complex of multiplication by that variable,
which is injective on the preceding polynomial
quotient.  Its degree-zero homology is the quotient
by the variables in that finite set.  Filtered
unions are exact for modules, so passage through
all finite initial sets gives positive homology
zero and degree-zero quotient $\mathbb Q$.

After tensoring with $\mathbb Q[t]$, we therefore
obtain the Koszul complex $K(a_N:N\geq N_0)$.
Write its degree-one generator as $\xi_N$, with
weight $N$ and $d\xi_N=t^N/N!$.  For $N>N_0$ put
\[
 \eta_N=\xi_N-
              \frac{N_0!}{N!}t^{N-N_0}\xi_{N_0}.
                                                               \tag{136.38}
\]
This is a homogeneous invertible change of basis
of the free direct sum on the generators: both
it and its inverse use only two terms for each
generator.  Moreover $d\eta_N=0$.  It follows that
\[
 K(a_N:N\geq N_0)
  \simeq K(t^{N_0}/N_0!)\otimes_{\mathbb Q[t]}
                 \Lambda_{\mathbb Q[t]}(\eta_N:N>N_0),
                  \qquad d\eta_N=0.                       \tag{136.39}
\]
Multiplication by $t^{N_0}/N_0!$ is injective in
$\mathbb Q[t]$.  The homology of its two-term complex
is therefore $\mathbb Q[t]/(t^{N_0})$ in degree zero.
The second factor in (136.39) is free in every
degree with zero differential, whence
\[
 H_*\bigl(K(a_N:N\geq N_0)\bigr)
  =\mathbb Q[t]/(t^{N_0})\otimes_{\mathbb Q}
                         \Lambda_{\mathbb Q}(\eta_N:N>N_0).
                                                               \tag{136.40}
\]
In particular $1\otimes\eta_N$ is a nonzero
horizontal $H_1$ class of weight $N$ for every
$N>N_0$.  The zero-form row is not horizontally
acyclic in any eventual range of weights, in
contradiction to condition (136.4).  Since flat
rationalization preserves homology, integral
eventual acyclicity would imply rational eventual
acyclicity.  The integral construction fails too.
\end{proof}

The argument applies even when $N_0=1$.  The
obstruction consists of the higher homology of
the redundant, independent coefficient relations;
it is distinct from the prime-weight $H_0$
obstruction in AQ-136-INTEGRAL-TAIL.  It refutes
this free independent-disk construction only.
Adding coherent higher relations or constructing
a different admissible source remains a separate
problem; the integral original conjecture (a)
and integral terminal localization are not
refuted by (136.40).
